% =========================================================================
% 三、模型假设与符号说明
% =========================================================================

\section{模型假设与符号说明}

以下把假设按\textbf{题面给定}、\textbf{解释与近似}、\textbf{算法与评价设计}三类分列。区分的意义是：第一类是题目条件，不作为本文的自设假设；第二类是可被质疑的口径，逐条给出依据与影响；第三类是实现方案的选择，不是客观事实，其取值在正文相应位置说明。

\subsection{题面给定的参数与条件}

\begin{enumerate}
  \item 储能最大容量 12000 kWh，允许储电量范围 $[1200,10800]$ kWh，最大充放电功率 5000 kW，充放电效率两向各 90\%（附录 1）；
  \item 允许供能余量（富余供能可弃置），微网不向外部电网售电、不产生售电收入；
  \item 问题一要求储能设备在 0:00 与 24:00 的储电量相同；
  \item 问题二、问题三、问题四中，供电低于负载时须按交易时刻电价的 5 倍\textbf{紧急购电}；
  \item 问题三、问题四的 Q3 分支中，计划购电量高于调整购电量的部分按交易时刻电价的 50\% 结算，调整购电量高于计划购电量的部分、其超出部分按交易时刻电价的 1.5 倍结算；
  \item 问题四的电价为附件 4 给出的逐 10 分钟区间实时电价，计划购电与调整购电均按该区间实时价结算。
\end{enumerate}

\subsection{解释与近似}

以下五条不是题面给定的事实，而是本文对题面未明示之处的读法。逐条给出依据与影响。

\begin{enumerate}
  \item \textbf{问题二的信息口径}：第 $d$ 天 0:00 决策时只可用第 $1,\dots,d-1$ 天的实测与附件 1 的曲线，当日及以后的实测值及其派生量一律不可用；决策规则超参数的选取同样受此约束。依据：题面要求``每天 0:00 制定当天计划''，计划不可能包含当日实测。影响：这是问题二与问题三全部预测逻辑的前提。
  \item \textbf{问题四的价格揭示口径}：附件 4 给出的是各区间\emph{已发生}的实时电价，故价格随区间推进逐步揭示——0:00 只知前一日及以前的价格，6:00、12:00、18:00 另知当日已过去区间的价格，任何决策均不使用尚未揭示区间的电价。影响：该口径的替代方案（日前已知全天价）在同一政策下相差 7.70 万元，其量化见 \S\ref{sec:q4-sens}。
  \item \textbf{功率数据按区间平均功率读}：附件 1、附件 2 的时间标记为区间末端，但题面未给出功率数据的定义。本文统一把附件值读作该 10 分钟\emph{区间}的平均功率，并以区间电量 $=P\,\Delta t$ 折算。影响：若实际定义为瞬时值，量纲折算需改为数值积分，本文的逐区间口径随之改变。
  \item \textbf{附件 3 预报的时间约定}：$f_\tau(h)$ 定义为\textbf{发布后第 $h$ 小时}的整点光伏功率预报，区间电量按相邻整点线性插值后在区间上积分得到（式~\eqref{eq:q3-pv}）。发布时刻 $\tau>0$ 时，首端点的功率取``上一已结束区间的平均功率''（由该区间电量除以 $\Delta t$ 换回 kW），这是对发布时刻瞬时功率的\textbf{端点近似}，不是直接观测到的瞬时值。影响：图~\ref{fig:q3-pv} 的对照表明该近似优于改用历史预报端点。
  \item \textbf{未建模的因素}：不引入题目未要求的电池退化成本、运维成本，不建模售电，也不把预测残差的时段间相关性显式写入规划模型。影响：这些因素会使实际的可行域与费用账本发生变化，故本文的全部结论限于上述口径。
\end{enumerate}

\subsection{算法与评价设计}

\begin{enumerate}
  \item \textbf{计划端的末端条件分问设定}：问题一的日循环是题面条件，直接写在模型约束中；问题二、问题三的计划层以``计划末端电量等于日初电量''作为单日规划的自洽边界，它\textbf{只约束计划轨迹，不约束实际运行}（实际日末电量可以不同于日初并结转到次日）；问题四的计划层把末端条件改为 48 小时时域末端回到起点（\S\ref{sec:q4-horizon}）。三者是不同问题的不同设计，不互相推广。
  \item \textbf{评价期与共同预热}：为消除各方案 1 月冷启动轨迹差异的干扰，所有可实施方案共用同一条 1 月预热轨迹，2 月 1 日 0:00 统一从 10800 kWh 起评，评价期为 2025.2.1--12.31 共 334 天；1 月仅用于提供历史数据与首次选参，不计入评价。
  \item \textbf{因果选参}：问题二的 $(\beta,W,\gamma,N_{\mathrm{DOW}})$ 按逐月 walk-forward 协议选取（\S\ref{sec:q2-walkforward}）；问题三的收缩系数 $\kappa_g$ 只用 $d$ 之前 30 天逐日重估。问题四的 $(\lambda_{\mathrm{p}},\delta)$ 与 48 小时时域为\textbf{同年探索性选定}——其逐时动作不读取未来信息，但候选范围的确定过程看过全年数据，两者不是一回事，本文在 \S\ref{sec:q4-selector}、\S\ref{sec:q4-limit} 如实标注，不宣称独立测试。
  \item \textbf{候选网格}：$W\in\{7,14\}$，$N_{\mathrm{DOW}}\in\{4,6,10\}$，$\gamma\in\{0,0.5,1.0\}$，$\beta\in\{0.70,0.75\}$（\S\ref{sec:q2-walkforward}）；问题四的 $\lambda_{\mathrm{p}}$、$\delta$ 按 \S\ref{sec:q4-grid} 的格点取定。网格边界处的最优性不作一般性断言。
  \item \textbf{分位裕量的独立性}：裕量按各时段单独取分位数，未建模时段之间的相关性；这将使同时段多区间同时缺口的联合概率不被显式控制。
\end{enumerate}

\subsection{时间离散口径}

附件 1/2 的时间标记为区间末端；区间 $t$ 指 $[t\Delta t,(t+1)\Delta t)$，$\Delta t = 1/6$ h，每日 $T=144$ 个区间。结合 \S 上文解释与近似中``功率数据按区间平均功率读''的口径，功率（kW）乘 $\Delta t$ 折算为电量（kWh）；电价为强度量，直接取附件值。问题一、二的结果文件均按官方模板原始行序填写，模板行标签较字面区间整体后移一个区间，核对时需注意。

\subsection{术语约定}

为保持全文用词一致，先作如下约定：

\begin{itemize}
  \item \textbf{预报}专指附件 3 给出的逐时光伏发电功率；\textbf{预测}专指本文模型给出的净负载估计；\textbf{点预测}指不含风险裕量的中心预测 $\mu$，\textbf{规划需求代理}指叠加裕量后进入规划的量 $\widetilde N$；
  \item \textbf{计划购电量}（记 $b^{0}$）指 0:00 制定的购电量；\textbf{最终合同购电量}（记 $b^{\mathrm f}$）指经调整层重定后实际执行与结算的购电量；\textbf{紧急购电量}指实测缺口超出储能释放能力后按 5 倍电价购入的电量；
  \item \textbf{供能余量}指已购或已发而未被消纳、最终弃置的电量；\textbf{光伏富余}指光伏出力高于同期负载的部分；
  \item \textbf{计划库存}（记 $\widehat E$）指线性规划轨迹上的储能电量；\textbf{实际库存}（记 $E$）指执行层按实测递推得到的储能电量；
  \item \textbf{实时电价}（仅见于问题四）指附件 4 给出的逐区间电价；\textbf{价格水平}指式~\eqref{eq:q4-decomp} 中的日水平因子 $\ell_d$；\textbf{价格形状}指归一化后的日内价格轮廓；
  \item \textbf{滚动时域}（问题四）指把计划线性规划的优化终点由当日 24:00 延至次日 24:00 的做法；\textbf{持久性预测}指以当日预测序列平铺来估计次日输入的近似方式；
  \item \textbf{方案}指一组可完整执行的配置（A、B、E、A0、A0K、B0、B1、F\_48、G\_48 等）；\textbf{策略}沿用题面含义，指逐日制定购电决策的规则。
\end{itemize}

\subsection{符号说明}

\begin{table}[htbp]
  \centering
  \caption{主要符号说明（问题一至问题三）}
  \label{tab:sym-common}
  \small
  \begin{tabularx}{0.9\textwidth}{cX}
    \toprule
    符号 & 含义 \\
    \midrule
    $t$ & 日内 10 分钟区间编号，$t=0,1,\dots,143$ \\
    $T=144$ & 每日区间数（固定）；$\Delta t=1/6$ h 为区间长度 \\
    $H$ & 单次优化问题的时域长度（区间数）：日计划 $H=T$，问题四的 48 小时计划 $H=2T$ \\
    $d$ & 日期编号，$d=0,\dots,364$（2025 年） \\
    $P^L_t,\ P^{PV}_t$ & 区间 $t$ 的负载、光伏平均功率（kW） \\
    $L_t,\ V_t$ & 区间 $t$ 的负载电量、光伏电量，$L_t=P^L_t\Delta t$，$V_t=P^{PV}_t\Delta t$（kWh） \\
    $N_t^{(d)}=L_t^{(d)}-V_t^{(d)}$ & 第 $d$ 天区间 $t$ 的\textbf{实测}净负载（kWh） \\
    $\mu_t$ & 净负载\textbf{点预测}（不含风险裕量，仅用历史信息） \\
    $\widetilde N_t$ & \textbf{规划需求代理}：点预测叠加风险裕量后进入线性规划的需求（kWh） \\
    $\widehat E_t,\ E_t$ & \textbf{计划}库存、\textbf{实际}库存（区间 $t$ 起始，kWh） \\
    $p_t$ & 区间 $t$ 的购电电价（元/kWh）；紧急电价为 $5p_t$ \\
    $b^{0}_t,\ b^{\mathrm f}_t$ & 0:00 计划购电量（日前合同）与最终合同购电量（kWh）；问题二无调整层，$b^{\mathrm f}_t=b^{0}_t$，正文简记 $b_t$ \\
    $c_t,\ d_t$ & 区间 $t$ 储能充电量、放电量（交流侧，kWh）；$\bar C=\bar D=5000/6$ kWh \\
    $e_t,\ w_t$ & 区间 $t$ 的紧急购电量、弃置供能余量（kWh） \\
    $\eta_c,\ \eta_d$ & 充电、放电效率，均为 $0.9$ \\
    $W,\ N_{\mathrm{DOW}}$ & 近期水平窗口长度（天）、同星期偏差所取的历史同星期日个数（天） \\
    $\gamma,\ \beta$ & 星期偏差校正强度、分位数裕量水平 \\
    \bottomrule
  \end{tabularx}
\end{table}

\begin{table}[htbp]
  \centering
  \caption{主要符号说明（问题三的预报通道与问题四的实时电价）}
  \label{tab:sym-q34}
  \small
  \begin{tabularx}{0.9\textwidth}{cX}
    \toprule
    符号 & 含义 \\
    \midrule
    $\tau$ & 预报发布时刻，$\tau\in\{0,6,12,18\}$（h）；对应区间 $t_\tau=6\tau$ \\
    $f_\tau(h)$ & 发布时刻 $\tau$ 对\textbf{发布后第 $h$ 小时}的整点光伏功率预报（kW）；$\widetilde V_t(\tau)$ 为其区间电量 \\
    $\kappa_g$ & 附件 3 预报通道的分档收缩系数，按预报时长分三档 \\
    $u_t,\ v_t$ & 合同相对 $b^{0}$ 的下调量、上调量（kWh），违约／加价分段 \\
    $\lambda$ & 门控实验中单决策点选用的预报融合强度，$\lambda\in\{0,\tfrac14,\tfrac12,1\}$ \\
    \midrule
    $P_{d,t}$ & 附件 4 给出的第 $d$ 天区间 $t$ 的实时电价（元/kWh）；结算时记 $p_t$ \\
    $\pi_t^{\mathrm{TOU}},\ \ell_d,\ r_{d,t}$ & 分时电价形状、日价格水平因子、价格残差（式~\eqref{eq:q4-decomp}） \\
    $\omega_t$ & 价格加权风险裕量的权重，由决策价形状按幂次生成，整日冻结 \\
    $\lambda_{\mathrm{p}},\ \delta$ & 价格加权强度（幂指数）与裕量水平偏移（式~\eqref{eq:q4-prm}） \\
    \bottomrule
  \end{tabularx}
\end{table}
